% !TeX root = main.tex
% Заполните значения ниже — остальные файлы менять не обязательно.

\newcommand{\ministry}{Министерство науки и высшего образования Российской Федерации}
\newcommand{\university}{ФЕДЕРАЛЬНОЕ ГОСУДАРСТВЕННОЕ АВТОНОМНОЕ\\ОБРАЗОВАТЕЛЬНОЕ УЧРЕЖДЕНИЕ ВЫСШЕГО ОБРАЗОВАНИЯ\\«НАЦИОНАЛЬНЫЙ ИССЛЕДОВАТЕЛЬСКИЙ УНИВЕРСИТЕТ ИТМО»\\(УНИВЕРСИТЕТ ИТМО)}
\newcommand{\faculty}{Факультет систем управления и робототехники}

\newcommand{\subject}{Адаптивное и робастное управление}
\newcommand{\worktype}{Отчёт по лабораторной работе}
\newcommand{\labnumber}{№3}
\newcommand{\variant}{1}

\newcommand{\studentlabel}{Работу выполнил:} % Для нескольких студентов: «Работу выполнили:»
\newcommand{\students}{Сухов Н. М.\\Иванов Н. М.}
% Несколько студентов: \renewcommand{\students}{Иванов И. И.\\Петров П. П.}
\newcommand{\group}{АдиРобУп (ИР) R23 бак 1.3.1}
\newcommand{\teacher}{Парамонов Алексей Владимирович}

\newcommand{\city}{Санкт-Петербург}
\newcommand{\reportyear}{\the\year}

% Метаданные итогового PDF.
\hypersetup{
  pdftitle={\subject. \worktype~№\labnumber},
  pdfauthor={\students},
  pdfsubject={\subject}
}
